% ===================================================================
%  BẢNG Số 2 — Cơ thể người. --- Giờ ra chơi. Các trò chơi.
%
%  Traduction, faite en 2026, en vietnamien standard d'aujourd'hui,
%  sur l'ido de texto/io. Regles de la colonne : voir 10-bang-01.tex.
%
%  LES TROIS COUCHES DU TABLEAU 1 SE RETROUVENT ICI, ET LA LIGNE PASSE
%  AU MEME ENDROIT. Le corps humain est vietnamien de part en part —
%  đầu la tete, mắt l'oeil, mũi le nez, tim le coeur, phổi le poumon,
%  xương l'os, máu le sang : pas un emprunt dans les quatre-vingts
%  mots de l'anatomie, parce que ces mots-la etaient dans la langue
%  avant toute ecole. Mais LES CINQ SENS sont sino-vietnamiens —
%  thị giác, thính giác, khứu giác, vị giác, xúc giác — et il le faut,
%  car ce ne sont pas des choses qu'on montre du doigt : ce sont des
%  notions de manuel, et le manuel parlait chinois.
%
%  ET LES JEUX SE PARTAGENT DE MEME. Ceux qui existaient ici gardent
%  leur nom vietnamien : bịt mắt bắt dê pour le colin-maillard —
%  litteralement « les yeux bandes, attraper la chevre » —, trốn tìm
%  pour cache-cache, cà kheo pour les echasses, con quay pour la
%  toupie, thả diều pour le cerf-volant. Ceux qui sont arrives dans
%  une caisse gardent leur nom francais, coupe en syllabes :
%  crô-kê, bi-bô-kê, quần vợt pour le tennis.
%
%  DEUX JEUX N'ONT PAS D'EQUIVALENT ET SONT DECRITS, NON NOMMES. Le
%  « jeu de barres » (38) et « l'ourse » (32) sont deux jeux d'equipe
%  francais que le vietnamien ne connait pas ; on les rend par ce
%  qu'ils font — trò cướp cờ, la prise du drapeau, et trò con gấu, le
%  jeu de l'ours. Le second garde son animal, que la note du livret
%  rend necessaire : le narrateur dit qu'il est bien francais.
%
%  « SOCCER » DU QUEBEC N'A PAS DE PROBLEME ICI : le vietnamien dit
%  bóng đá, la balle au pied, et le mot ne se confond avec rien.
% ===================================================================

%%K t02-apar-1 apar
\VUsaut{4.0mm}
\VUtitre{15.0pt}{60}{BẢNG Số 2}
\VUsaut{2.0mm}
\VUtitre{11.4pt}{0}{Cơ thể người. --- Giờ ra chơi.}
\VUtitre{11.4pt}{0}{Các trò chơi.}

%%K t02-tit-0 sub
\VUtitre{13.2pt}{0}{Cơ thể người.}
\VUsaut{2.4mm}
\VUcentre{10.2pt}{0}{\textit{(Bài học đơn giản về khoa học tự nhiên.)}}
\VUsaut{3.0mm}

%%K t02-tit-1 sub pos=t02-01-1
\VUpk{10.2pt}{40}{Cái đầu.}
\VUsaut{3.0mm}

%%K t02-01-1 p
1. --- Những phần chính của cơ thể người là : \VUgras{đầu}, \VUgras{thân mình} và
\VUgras{tay chân}.

%%K t02-02-1 p
2. --- Đầu gồm hai phần : \VUgras{sọ}, chứa \VUgras{não} và được \VUgras{tóc} che
phủ, và \VUgras{khuôn mặt}.

%%K t02-03-1 p
3. --- Trên hình vẽ của chúng ta không thấy sọ cũng không thấy não; nhưng ta thấy rất
rõ những phần quan trọng của khuôn mặt.

%%K t02-04-1 p
4. --- Trước hết là \VUgras{trán}, cho thấy một trí tuệ mạnh mẽ, một tư duy luôn hoạt
động. \VUgras{Thái dương} rất thoáng. \VUgras{Mắt} linh lợi và mở to. \VUgras{Lông
mày} rậm và \VUgras{lông mi} dài che chở cho mắt.

%%K t02-05-1 p
5. --- Tiếp đến là \VUgras{mũi} và \VUgras{lỗ mũi}. Mũi giúp ta ngửi và thở. Hai bên
mũi là \VUgras{má}, xa hơn nữa là \VUgras{tai}. Phần dưới của khuôn mặt gồm
\VUgras{miệng} và \VUgras{cằm}.

%%K t02-06-1 p
6. --- Ta không thấy rõ lắm, vì \VUgras{ria mép} và \VUgras{tóc mai} che khuất chúng.
Nhưng ta biết rằng miệng có hai \VUgras{môi}, \VUgras{hàm trên} và \VUgras{hàm dưới}
với \VUgras{răng} của chúng, và \VUgras{lưỡi}. Nhờ miệng ta nói và ăn. Nhờ lưỡi ta nếm
được thức ăn.

%%K t02-07-1 p
7. --- Mắt, tai, mũi và miệng, cũng như các ngón tay, là cơ quan của năm giác quan :
\VUgras{thị giác}, \VUgras{thính giác}, \VUgras{khứu giác}, \VUgras{vị giác},
\VUgras{xúc giác}.

%%K t02-08-1 p
8. --- Vậy đầu là một trong những phần thú vị nhất của cơ thể người.

%%K t02-tit-2 sub
\VUsaut{4.4mm}
\VUpk{10.2pt}{40}{Thân mình.}
\VUsaut{3.0mm}

%%K t02-09-1 p
9. --- \VUgras{Cổ} nối đầu với thân mình. Nó gồm \VUgras{cổ họng} và \VUgras{gáy}.

%%K t02-10-1 p
10. --- Thân mình cũng chứa nhiều cơ quan thiết yếu. Trong \VUgras{lồng ngực} có
\VUgras{phổi}, nhờ đó ta thở, và \VUgras{tim}, trung tâm của sự sống. Khi tim ngừng
đập, sinh vật chết.

%%K t02-11-1 p
11. --- \VUgras{Xương sườn} giữ lấy lồng ngực.

%%K t02-12-1 p
12. --- Những phần khác của thân mình là \VUgras{hông}, \VUgras{lưng}, \VUgras{vai} và
\VUgras{bụng}. Ta nằm nghiêng hoặc nằm ngửa để ngủ; ta vác đồ nặng trên vai.

%%K t02-13-1 p
13. --- Trong bụng có \VUgras{dạ dày} và \VUgras{ruột}. Chúng giúp ta tiêu hóa thức ăn.

%%K t02-tit-3 sub
\VUsaut{4.4mm}
\VUpk{10.2pt}{40}{Tay chân.}
\VUsaut{3.0mm}

%%K t02-14-1 p
14. --- Ta có bốn \VUgras{chi} : hai \VUgras{cánh tay} và hai \VUgras{chân}. Nhờ cánh
tay ta cầm, giữ, nâng và nhặt đồ vật; nhờ chân ta đứng, đi và chạy.

%%K t02-15-1 p
15. --- \VUgras{Vai} nối cánh tay với thân mình. Một \VUgras{bắp thịt} rất khỏe, bắp
tay, làm cánh tay cử động; \VUgras{khuỷu tay} nối cánh tay với \VUgras{cẳng tay}.
\VUgras{Cổ tay} tạo thành khớp của \VUgras{bàn tay}, tận cùng bằng các \VUgras{ngón
tay} và \VUgras{móng tay}. Trên bàn tay ta phân biệt được \VUgras{tĩnh mạch} (và động
mạch), nơi \VUgras{máu} chảy.

%%K t02-16-1 p
16. --- Các phần của chân là : \VUgras{đùi}, \VUgras{đầu gối} và \VUgras{khoeo chân},
\VUgras{bắp chân}, mà \VUgras{thịt} được \VUgras{da} bọc lấy, \VUgras{mắt cá} và
\VUgras{bàn chân}. \VUgras{Xương} chân rất to và rất chắc. Ở đầu bàn chân là các
\VUgras{ngón chân}. Những phần còn lại là \VUgras{gan bàn chân} và \VUgras{gót chân}.

%%K t02-17-1 p
17. --- Đó là bản mô tả gần như đầy đủ về cơ thể người.

%%K t02-17-2 p
\textit{Chú thích.} --- \textit{Với cách nói « tôi đau ở... (đầu, v.v.) », thầy giáo có
thể dùng lại toàn bộ vốn từ này dưới góc độ} \VUgras{thể trạng} \textit{tốt hay xấu.}

%%K t02-tit-4 sub
\VUsaut{5.0mm}
\VUpk{10.2pt}{40}{Thể dục.}
\VUsaut{2.4mm}
\VUcentre{10.2pt}{0}{\textit{(Chuyện kể của một học sinh.)}}
\VUsaut{3.0mm}

%%K t02-18-1 p
18. --- Để được mềm dẻo, nhanh nhẹn và khỏe mạnh, chúng em phải luyện chân và tay. Vì
thế chúng em chơi và tập \VUgras{thể dục}.

%%K t02-19-1 p
19. --- Trong tuần, hai việc ấy diễn ra ở sân trường trung học (xem bảng số 1). Nhưng
thứ năm và chủ nhật, chúng em ra một bãi cỏ rộng, nơi có chỗ để chạy và chơi đùa.

%%K t02-20-1 p
20. --- Thầy cô và cha mẹ chúng em đôi khi cùng đi; nhưng chính \VUgras{thầy dạy thể
dục}\textsuperscript{(12)} là người hướng dẫn các bài tập.

%%K t02-21-1 p
21. --- Trên bãi cỏ, bên trái lối vào, có \VUgras{dụng cụ thể
dục}\textsuperscript{(1)}; chúng em trèo \VUgras{thang}\textsuperscript{(2)}, lộn người
trên \VUgras{vòng treo}\textsuperscript{(3)} và trên \VUgras{đu bay}\textsuperscript{(4)},
đu bằng tay lên tận đỉnh \VUgras{thang dây}\textsuperscript{(5)} hoặc \VUgras{dây thừng
có nút}\textsuperscript{(6)}.

%%K t02-22-1 p
22. --- Trong lúc đó, các bạn em nhảy qua \VUgras{ngựa gỗ}\textsuperscript{(7)} hoặc
\VUgras{xà kép}\textsuperscript{(11)}. Những bạn khác thì \VUgras{đấm
bốc}\textsuperscript{(8)} hoặc \VUgras{đấu kiếm}\textsuperscript{(9)} với mặt nạ và
\VUgras{kiếm nhẹ}. Sau cùng, có bạn \VUgras{nâng tạ}\textsuperscript{(10)}.

%%K t02-tit-5 sub
\VUsaut{4.6mm}
\VUpk{10.2pt}{40}{Các trò chơi.}
\VUsaut{3.0mm}

%%K t02-23-1 p
23. --- Chung quanh những bạn tập thể dục, các em trai và em gái chơi nhiều \VUgras{trò
chơi} khác nhau. Các em gái chơi \VUgras{vòng}\textsuperscript{(14)}; các em nhảy
\VUgras{dây}\textsuperscript{(13)} hoặc rủ anh em và anh chị em họ chơi một ván
\VUgras{crô-kê}\textsuperscript{(15)}. Các em thích thú đánh bật \VUgras{quả bóng} của
đối phương bằng \VUgras{vồ gỗ} của mình, đúng lúc bạn ấy tưởng sẽ dễ dàng qua được cổng
vòng!

%%K t02-24-1 p
24. --- Các em trai thích những trò dùng sức hơn. Các em chơi
\VUgras{ki}\textsuperscript{(16)}, đánh đổ chúng rất khéo bằng \VUgras{quả bóng}\textsuperscript{(17)}. Các em cũng không
chê \VUgras{con quay}\textsuperscript{(18)} và \VUgras{bi-bô-kê}\textsuperscript{(19)},
hay cả \VUgras{trò ném ếch}\textsuperscript{(20)}. Nhưng trò các em thích nhất vẫn là
\VUgras{bi}\textsuperscript{(21)} và \VUgras{bóng ném tường}\textsuperscript{(22)}, vì
bức tường của \VUgras{nhà kính}\textsuperscript{(26)} rất thuận tiện để quả bóng nhẹ nảy
lại.

%%K t02-25-1 p
25. --- Bé Gioan, đi trên \VUgras{cà kheo}\textsuperscript{(24)}, rất khéo và giữ thăng
bằng rất giỏi. Gần em, mấy người bạn chơi \VUgras{trốn tìm}\textsuperscript{(25)} trong
nhà kính và sau \VUgras{hàng rào}. Những bạn khác thích \VUgras{trò đập
tay}\textsuperscript{(28)} hoặc \VUgras{quả cầu}\textsuperscript{(29)} bay từ
\VUgras{vợt} này sang vợt kia.

%%K t02-noto-1 noto
\VUnotes{143.2mm}{%
(*) Trò chơi của trẻ con thường mang những cái tên thuần túy dân tộc. Chúng tôi đã cố
gắng đặt tên cho chúng bằng tiếng ido một cách tốt nhất có thể, hoặc dựa vào chính hành
động đặc trưng của trò chơi, hoặc chọn tên dân tộc của nước này nước kia, khi tên ấy
không quá sai lệch. Ví dụ : \VUgras{trò thùng phuy} (tiếng Đức và tiếng Pháp) chính là
\VUgras{trò ném ếch} \textsuperscript{(20)} (tiếng Anh), trong đó người chơi cố ném
những cái đĩa tròn của mình vào miệng con ếch bằng kim loại đang há ra trên cái thùng
(cái phuy) đục lỗ.
\VUgras{Trò nhảy vai}\textsuperscript{(30)} chỉ khác \VUgras{trò nhảy cừu} ở điểm này :
để nhảy qua đầu, người chơi chống tay lên vai bạn mình, còn trong « \VUgras{trò nhảy
cừu} » thì chỉ chống tay lên lưng. --- Hoặc nữa, vài người tham gia cúi xuống trong khi
những người khác nhảy qua lưng họ, chống tay lên vai; nhóm thứ nhất phải chịu được cú va
mà không khuỵu xuống, nhóm thứ hai phải nhảy mà không mất thăng bằng (xem chính bảng
này).%
}

%%K t02-26-1 p
26. --- Giữa bãi cỏ có hai cái cây. Dưới gốc một trong hai cây ấy, bảy tám em trai đã
tổ chức một ván \VUgras{nhảy vai}\textsuperscript{(30)} (*). Trò này không phải nước nào
cũng biết; nhưng ai cũng biết \VUgras{trò nhảy cừu} là gì, dù lần này các em không chơi
trò ấy.

%%K t02-tit-6 sub
\VUsaut{5.0mm}
\VUpk{10.2pt}{40}{Trò chơi ngoài trời.}
\VUsaut{3.4mm}

%%K t02-27-1 p
27. \VUgras{Bịt mắt bắt dê}\textsuperscript{(31)} cũng rất vui; các em gái và em trai
lần lượt đeo \VUgras{khăn bịt mắt} lên mắt mình và bắt những bạn vụng về để mình bị tóm.

%%K t02-28-1 p
28. --- Giữa bãi cỏ, đây là một trò rất Pháp nhưng hơi thô bạo : đó là \VUgras{trò con
gấu}\textsuperscript{(32)}, ồn ào hơn nhiều so với trò \VUgras{thả
diều}\textsuperscript{(33)} êm ả, hay so với trò của người Anh là \VUgras{quần
vợt}\textsuperscript{(34)}.

%%K t02-29-1 p
29. --- Môn thể thao sau cùng ấy là niềm vui của các anh lớn. Chính các thầy giáo cũng
sẵn lòng chơi. Nó đòi hỏi nhiều khéo léo và nhanh nhẹn, và em rất thích. Còn thú vui của
\VUgras{xích đu}\textsuperscript{(35)} thì em nhường cho các em gái.

%%K t02-30-1 p
30. --- Em không thích một trò khác của người Anh, một trò có thể trở nên nguy hiểm.

%%K t02-30-1b p
Em muốn nói đến \VUgras{bóng đá}\textsuperscript{(36)}, môn làm anh cả của em vui sướng.
Em thích hơn \VUgras{trò cướp cờ}\textsuperscript{(38)} xưa cũ của chúng ta, cũng tốt
cho sức khỏe như thế mà thật sự ít thô bạo hơn.

%%K t02-tit-7 sub
\VUsaut{4.6mm}
\VUpk{10.2pt}{40}{Xe đạp và trò chơi bóng.}
\VUsaut{3.0mm}

%%K t02-31-1 p
31. --- Em cũng thích đi một vòng quanh bãi cỏ bằng \VUgras{xe đạp}\textsuperscript{(37)}.

%%K t02-32-1 p
32. --- Sang năm em sẽ học \VUgras{cưỡi ngựa}\textsuperscript{(39)}. Con ngựa đi hay
chạy nước kiệu một cách duyên dáng, rồi phi nước đại nhảy qua chướng ngại vật, đẹp biết
bao!

%%K t02-33-1 p
33. --- Mùa hè, chúng em học \VUgras{bơi}\textsuperscript{(40)}. Chúng em thích thú lặn
và tắm trong làn nước trong mát của dòng sông. Đôi khi, sau cùng, chúng em thả một
\VUgras{khinh khí cầu} bằng giấy \textsuperscript{(41)}, nó bay lên trời một cách oai vệ
ngay khi được bơm đầy khí nóng.

%%K t02-34-1 p
34. --- Còn nhiều trò chơi khác nữa, nhưng em kể mãi cũng không hết, và qua bản tóm tắt
này ai cũng hiểu rằng vào những ngày thứ năm, người ta chẳng chán chút nào trên bãi cỏ
đẹp đẽ của chúng em.

%%K t02-34-2 p
T.B. --- \textit{Thầy giáo sẽ dễ dàng bổ sung chương này bằng cách mô tả tỉ mỉ hơn từng
trò chơi quan trọng nhất.}
\VUsaut{6.0mm}
\VUfilet{22mm}
